%% !TEX root = manuscript.tex

\chapter{Contexte et enjeux}
\minitoc
\newpage

Ce premier chapitre pose le cadre de mon alternance. Après une présentation de l'entreprise et
de mon rattachement, j'y décris le système d'information du groupe et les données que j'ai
manipulées, avant d'introduire la problématique qui guide l'ensemble de ce mémoire. Je termine
en présentant brièvement les outils qui reviennent tout au long du document.

\section{Contexte métier : le groupe Angelotti et la Holding SPA}

J'ai effectué mon alternance de Master 2 au sein de la Holding SPA du groupe Angelotti, un
acteur de l'aménagement et de la promotion immobilière implanté en Occitanie et en région PACA.

Ces deux métiers n'ont pas la même logique économique. Dans l'aménagement, la valeur naît de la
transformation du foncier, c'est-à-dire de la division d'un terrain en lots viabilisés. Dans la
promotion, elle naît de la construction de logements neufs. Ils produisent donc deux structures
de coûts et deux rythmes de commercialisation différents, ce qui explique que je les traite
séparément tout au long de ce mémoire.

Une deuxième particularité du secteur a pesé sur mon travail. Chaque opération étant portée par
une société dédiée, le suivi financier est nativement cloisonné : il n'existe pas de vue
d'ensemble spontanée sur l'activité. Une opération vit par ailleurs sur un cycle qui se compte
en années. Un budget est arrêté au moment de l'engagement, puis le chantier et la
commercialisation le font dériver, parce que les appels d'offres reviennent plus chers
qu'anticipé, que des acquéreurs se désistent ou que le rythme de vente ralentit. Cette dérive
est le fonctionnement normal du métier. Ce qui pose problème, c'est de la mesurer à temps, car
sur une marge de quelques centaines de milliers d'euros, quelques mois de retard suffisent à
transformer un ajustement en perte.

La Holding SPA regroupe les fonctions support du groupe. Rattachée à elle, j'ai travaillé sous
l'encadrement de ma tutrice, Sabine Cruanas, et en lien direct avec les services métiers : la
direction financière, le service juridique et les équipes commerciales. Le groupe ne dispose
pas d'une équipe data constituée, ce qui explique l'étendue du périmètre qui m'a été confié :
j'y ai été, sur l'année, la seule personne dédiée à ces sujets.

\section{Contexte data : le système d'information, les données et leur refonte}

\subsection{Le système d'information et sa refonte}

À mon arrivée en deuxième année, le système d'information reposait sur trois outils principaux.
\textbf{Grimmo}, l'ERP métier historique, pour la gestion des opérations ; \textbf{Vadimm}, le
CRM, pour la relation commerciale ; et \textbf{MyReport}, l'outil décisionnel, pour les tableaux
de bord des services. À ces trois outils s'ajoutait une quatrième couche, moins visible :
\textbf{les tableurs}, qui faisaient le lien partout où les trois premiers ne se parlaient pas.

La figure~\ref{fig:si} représente cette organisation. Les trois outils et leurs usages métier
occupent le haut du schéma, reliés par les flèches qui figurent la circulation normale de la
donnée. Les tableurs forment la bande du bas, et les traits pointillés qui les rattachent à
chacun d'eux marquent les échanges qui ne passent par aucun système. On y lit le problème que
la reprise des données allait devoir traiter : le même client pouvait exister dans le CRM et
dans l'ERP sans qu'aucun lien ne les rattache.

\begin{figure}[!htbp]
\centering
\includegraphics[width=0.65\textwidth]{fig_1_1_systeme_information.png}
\caption{Les outils du système d'information et la circulation de la donnée}
\label{fig:si}
\end{figure}

C'est dans ce décor que s'est engagée la refonte : le remplacement de Grimmo par SPO, la
solution éditée par Salvia~[2]. Ce chantier a commandé le calendrier de toute mon année, et
trois contraintes en découlaient. Une contrainte de \textbf{volumétrie} d'abord, puisque 283
opérations devaient être reprises avec leurs subdivisions et leurs budgets, échelle à laquelle
une saisie manuelle intégrale était exclue. Une contrainte de \textbf{continuité de service}
ensuite, les utilisateurs devant pouvoir facturer dès le basculement. Une contrainte de
\textbf{conformité au modèle cible} enfin, SPO imposant ses propres règles d'intégrité, qui
n'étaient pas toutes documentées et qu'il a fallu découvrir en les heurtant.

À ces contraintes s'est ajoutée une perspective absente du calendrier initial, le passage vers
l'infrastructure technique du groupe Nexity, qui a influé sur certains choix techniques.

\subsection{Les données disponibles}

Les données que j'ai manipulées sont de deux natures. D'un côté, les données de structure et de
gestion : les opérations, leurs tranches de travaux et de commercialisation, et leurs budgets,
issus de Grimmo puis de SPO. De l'autre, les données commerciales : les clients acquéreurs,
leurs coordonnées, les lots, les réservations et les désistements, issus du CRM Vadimm. Le
projet du second semestre a par ailleurs mobilisé quatre exports du système de gestion et trois
sources publiques, que je détaille au chapitre~3.

Une partie de ces données est sensible au regard du RGPD, puisqu'il s'agit de données
personnelles de clients : identité, adresse postale, courriel et téléphone. La confidentialité
a donc été une contrainte permanente, à la fois dans mon travail et dans la rédaction de ce
mémoire. Les chiffres que je présente sont des volumétries ou des agrégats, et aucune donnée
nominative n'apparaît, ni dans ce document ni dans l'application que j'ai livrée.

\subsection{Mon positionnement au regard de la data science}

Je dois préciser d'emblée un point, parce qu'il éclaire tout ce qui suit : sur les quatre
missions de l'année, une seule relève de la data science au sens où le Master l'entend, et
c'est celle du second semestre. Le reste tient de l'ingénierie des données et de la gestion
d'applications.

Cette répartition ne m'a pas éloignée de la formation, elle m'a plutôt fait comprendre ce qui
la précède. En travaillant au plus près des données sources, j'ai vu concrètement pourquoi
certaines analyses sont possibles et d'autres non, et ce que coûte une donnée mal structurée en
amont. Ce travail de fond garantit une base saine pour les analyses à venir, et il alimente
directement l'outil de pilotage que j'ai construit ensuite.

\section{Problématique et fil conducteur}

Ces différents aspects se rejoignent autour d'un même constat. À chaque mission, avant de
pouvoir répondre à la question posée par un service, il m'a fallu commencer par reconstituer la
donnée : décider ce qu'est un client, ce qu'est une marge, ou ce que contient réellement un
fichier dont le nom annonce autre chose. Et à chaque fois, la qualité de cette reconstruction a
déterminé ce que je pouvais produire ensuite.

Qu'il s'agisse de migrer un référentiel vers un nouvel ERP, de maintenir un reporting ou de
construire un outil d'aide à la décision, l'objectif est resté le même : fournir aux services
des données justes, sur lesquelles ils peuvent réellement s'appuyer pour décider. C'est autour
de cet enjeu que s'articule ma problématique.

\begin{center}
\fbox{\begin{minipage}{0.88\textwidth}
\vspace{0.4em}
\begin{center}
\textbf{Comment reconstituer et fiabiliser les données d'une entreprise\\
au cours d'un changement d'ERP, afin d'en tirer des outils de pilotage\\
utilisables par les services métiers ?}
\end{center}
\vspace{0.4em}
\end{minipage}}
\end{center}

Cette question comporte deux volets, que mes deux missions principales illustrent chacune de
leur côté. Le premier est celui de la structure : reconstituer un socle de données dans un
modèle cible dont les règles ne sont pas documentées. Le second est celui de l'exploitation :
tirer de ce socle des analyses qui aident réellement à décider. Le second dépend du premier,
et c'est ce lien que je m'attache à montrer tout au long de ce mémoire.

\section{Structure du mémoire}

Après ce chapitre de cadrage, je présente les deux grandes missions de l'année. La
\textbf{Mission~1} (chapitre~2) porte sur la reprise des données vers SPO, conduite en deux
temps : les opérations au premier semestre, les tiers au second. La \textbf{Mission~2}
(chapitre~3) porte sur le Copilote Financier, le projet de Data Science mené pour la direction
financière. Je présente ensuite les \textbf{activités complémentaires} menées en parallèle sur
les deux semestres, avec le maintien du système décisionnel et le déploiement du SSO
(chapitre~4). Le chapitre~5 revient enfin sur mon organisation, les compétences acquises,
l'adéquation de ces missions avec le Master MIASHS, le positionnement RSE de l'entreprise et
mon projet professionnel.

\section{Les outils et l'environnement technique}

Mes missions s'appuient sur un ensemble d'outils qui structurent mon travail au quotidien.
Certains constituent la chaîne de traitement de la donnée elle-même, d'autres encadrent
l'organisation et la documentation. Comme ils reviennent tout au long de ce mémoire, je les
présente brièvement ici.

\textbf{SPO} (Suivi des Programmes Opérationnels) L'ERP métier qui remplace Grimmo depuis cette
année, édité par Salvia. Il porte l'ensemble de la gestion des opérations : structures,
budgets, tiers et facturation. C'est la cible de toutes les reprises de données que je décris
au chapitre~2.

\textbf{DirectAPI} Le module de l'éditeur Salvia qui permet d'injecter dans l'ERP des fichiers
tabulaires, sans passer par la saisie manuelle. C'est par lui que sont passés les gros volumes
structurels de la migration.

\textbf{L'API de SPO} L'interface de programmation de l'ERP, qui expose les mêmes objets que
l'interface graphique mais permet de les créer et de les modifier par requêtes. Je m'en suis
servie là où le format d'import tabulaire ne suffisait pas, en m'appuyant notamment sur le
format JSON Patch~[3].

\textbf{Postman} Un client d'interface de programmation, qui permet de composer des requêtes,
de les organiser en collections et de les exécuter par lots. C'est l'outil avec lequel j'ai
exécuté les 872 requêtes de modification décrites au chapitre~2.

\textbf{Vadimm} Le CRM du groupe, qui porte la relation commerciale : les clients, les
réservations et les désistements. C'est la source des données reprises au second semestre.

\textbf{MyReport} L'outil décisionnel du groupe, qui requête les bases sources et produit les
tableaux de bord des services. Son module \emph{Builder} sert à construire les extractions, et
c'est aussi par lui que passent les reportings que je maintiens.

\textbf{Le tableur} Excel occupe une place particulière dans ce mémoire. Il n'est pas seulement
un outil de saisie : je l'ai utilisé comme un véritable moteur de transformation, les fichiers
de travail que je décris au chapitre~2 étant entièrement pilotés par formules.

\textbf{Python} Le langage dans lequel j'ai écrit toute la chaîne du Copilote Financier :
nettoyage des exports, construction de la base, analyses et application. Les bibliothèques
mobilisées sont détaillées au chapitre~3.

\textbf{SQL et SQLite} SQL est le langage d'interrogation des bases relationnelles. SQLite est
le moteur léger dans lequel j'ai matérialisé la base analytique du Copilote Financier, avec ses
tables et ses vues.

\textbf{Streamlit} Une bibliothèque Python qui permet de construire une application web à partir
d'un script. C'est avec elle qu'est réalisée la plateforme livrée à la direction financière.

\textbf{Git et GitHub} Git est un système de gestion de versions, qui historise chaque
modification du code ; GitHub est la plateforme qui héberge le dépôt. C'est ce qui rend le
travail du second semestre traçable et reproductible.

\textbf{SSO} (Single Sign-On) L'authentification unique, qui donne accès à plusieurs
applications avec un seul compte. J'ai contribué à son déploiement sur deux périmètres du
groupe, comme je l'expose au chapitre~4.
